\documentclass[carta,final]{lafrox}
\usepackage{simulacion}
\input{datos-proponente}
\usepackage{needspace}% keep headings with their content
\let\SIMoldsection\section
\renewcommand{\section}{\needspace{6\baselineskip}\SIMoldsection}
\begin{document}
\SIMtitulo{Escritura de constitución y modificaciones}{Extracto de escritura pública --- LafroX SpA}

\noindent\small\textit{Extracto recreado con fines académicos para el Art.~39°, numeral~1, de las Bases Administrativas. Notaría, repertorios, inscripciones y personas son ficticios.}\normalsize

\section*{I. Escritura de constitución}
\noindent\textbf{REPERTORIO N.\textdegree{} SIM-2012-1.118.} \textbf{CONSTITUCIÓN DE SOCIEDAD POR ACCIONES LAFROX SpA.} En Valparaíso, a \SIMconstitucion, ante mí, \SIMnotarioTitular, Notario Público titular de la \SIMnotarioNombre, comparecen:

\begin{tablalafrox}{Comparecientes y accionistas fundadores.}{tab:esc1}{Y L{2.6cm} L{2.2cm} L{3.4cm}}
  {\cab{Compareciente} & \cab{RUN} & \cab{Nacionalidad} & \cab{Profesión}}
  \SIMsocioAnombre & \SIMsocioArun & chilena & ingeniero civil informático \\
  \SIMsocioBnombre & \SIMsocioBrun & chilena & ingeniera comercial \\
  \SIMsocioCnombre & \SIMsocioCrun & chilena & ingeniero civil industrial \\
  \SIMsocioDnombre & \SIMsocioDrun & chilena & ingeniera en información y control de gestión \\
\end{tablalafrox}

\noindent Los comparecientes, mayores de edad, acreditan su identidad con las cédulas citadas y exponen que constituyen una sociedad por acciones regida por la Ley N.\textdegree{} 18.046 y por las estipulaciones siguientes:

\begin{description}
  \item[PRIMERO. Nombre.] La sociedad se denominará \textbf{LafroX SpA}, pudiendo usar el nombre de fantasía \textbf{LafroX}.
  \item[SEGUNDO. Domicilio.] El domicilio de la sociedad será la comuna de Valparaíso, sin perjuicio de las agencias o sucursales que se establezcan en otros lugares del país o del extranjero.
  \item[TERCERO. Objeto.] La sociedad tendrá por objeto: a) el diseño, desarrollo, integración, comercialización y operación de sistemas y plataformas de software; b) la consultoría informática, la gestión de infraestructura tecnológica y la prestación de servicios de tecnologías de la información, incluida la mesa de servicios y el monitoreo de operaciones; c) la comercialización y soporte de equipamiento tecnológico; y d) la realización de toda otra actividad conexa o complementaria que acuerden los accionistas.
  \item[CUARTO. Duración.] La duración de la sociedad será indefinida.
  \item[QUINTO. Capital.] El capital es de \SIMcapitalinicial, dividido en 100.000 acciones nominativas, de una misma serie y sin valor nominal, suscritas y pagadas íntegramente por los accionistas en las proporciones de la Tabla~\ref{tab:esc2}.
  \item[SEXTO. Administración.] La sociedad será administrada por un Directorio de tres miembros y un Gerente General. Se designa Presidenta del Directorio a \SIMpresidenta, RUN \SIMpresidentarun. Se designa Gerente General a \SIMgerente, RUN \SIMgerenterun, quien será el representante legal de la sociedad y la representará judicial y extrajudicialmente con las facultades de un mandatario con libre administración y las que le otorgue el Directorio, pudiendo delegar parte de ellas en apoderados.
  \item[SÉPTIMO. Juntas.] Los accionistas se reunirán en junta ordinaria dentro del primer cuatrimestre de cada año y en junta extraordinaria cuando lo exijan las necesidades sociales. Los acuerdos se adoptarán por mayoría absoluta de las acciones emitidas con derecho a voto, salvo los casos que la ley o estos estatutos exijan un quórum superior.
  \item[OCTAVO. Balance y utilidades.] La sociedad confeccionará un balance general al 31 de diciembre de cada año. Salvo acuerdo unánime en contrario, se distribuirá como dividendo al menos el treinta por ciento de las utilidades líquidas.
  \item[NOVENO. Disolución y arbitraje.] La sociedad se disolverá por las causales legales. Las diferencias entre los accionistas, o entre éstos y la sociedad, se someterán a arbitraje de derecho ante árbitro designado de común acuerdo o, a falta de éste, por la justicia ordinaria de Valparaíso.
\end{description}

\begin{tablalafrox}{Suscripción y pago del capital inicial.}{tab:esc2}{Y R{2.6cm} R{2.4cm} R{2.4cm}}
  {\cab{Accionista} & \cab{Acciones} & \cab{Participación} & \cab{Aporte (CLP)}}
  \SIMsocioAnombre & 40.000 & 40 \% & 40.000.000 \\
  \SIMsocioBnombre & 25.000 & 25 \% & 25.000.000 \\
  \SIMsocioCnombre & 20.000 & 20 \% & 20.000.000 \\
  \SIMsocioDnombre & 15.000 & 15 \% & 15.000.000 \\
  \midrule
  \textbf{Total} & \textbf{100.000} & \textbf{100 \%} & \textbf{100.000.000} \\
\end{tablalafrox}

\noindent Un extracto de esta escritura se inscribió a fojas \SIMcbrfojas{} número \SIMcbrnumero{} del Registro de Comercio de \SIMcbrciudad{} del año \SIMcbranio{} y se publicó en el medio que la ley exige.

\section*{II. Primera modificación: aumento de capital (2018)}
\noindent\textbf{REPERTORIO N.\textdegree{} SIM-2018-3.402.} Por escritura pública de 22 de agosto de 2018, otorgada en la \SIMnotarioNombre, la junta extraordinaria de accionistas acordó aumentar el capital social de \SIMcapitalinicial{} a \SIMcapital, mediante la capitalización de utilidades retenidas por \$ 400.000.000 y el aporte en dinero de los accionistas, a prorrata, por \$ 500.000.000. Asimismo, se acordó una \textbf{división de acciones en razón de 1:10}, de modo que cada una de las 100.000 acciones existentes quedó dividida en diez acciones, sin alterar la participación de ningún accionista. En consecuencia, el artículo QUINTO de los estatutos queda así: «El capital de la sociedad es de \SIMcapital, dividido en \SIMaccionestotal{} de acciones nominativas, de una misma serie y sin valor nominal». Se inscribió al margen de la inscripción matriz, a fojas 3.502 número 1.644 del año 2018.

\section*{III. Segunda modificación: domicilio de operaciones (2019)}
\noindent\textbf{REPERTORIO N.\textdegree{} SIM-2019-2.775.} Por escritura pública de 10 de septiembre de 2019, otorgada en la \SIMnotarioNombre, la junta extraordinaria acordó reemplazar el artículo SEGUNDO de los estatutos por el siguiente: «El domicilio de la sociedad será la comuna de Valparaíso, y su centro de operaciones funcionará en Av.~Brasil 2241, Valparaíso, sin perjuicio de su sede central y de las agencias o sucursales que establezca en otros lugares». Se inscribió a fojas 5.117 número 2.286 del año 2019.

\begin{tablalafrox}{Estructura de capital vigente.}{tab:esc3}{Y R{2.6cm} R{2.4cm} R{3cm}}
  {\cab{Accionista} & \cab{Acciones} & \cab{Participación} & \cab{Capital (CLP)}}
  \SIMsocioAnombre & \SIMsocioAacc & \SIMsocioApct{} \% & 400.000.000 \\
  \SIMsocioBnombre & \SIMsocioBacc & \SIMsocioBpct{} \% & 250.000.000 \\
  \SIMsocioCnombre & \SIMsocioCacc & \SIMsocioCpct{} \% & 200.000.000 \\
  \SIMsocioDnombre & \SIMsocioDacc & \SIMsocioDpct{} \% & 150.000.000 \\
  \midrule
  \textbf{Total} & \textbf{\SIMaccionestotal} & \textbf{100 \%} & \textbf{1.000.000.000} \\
\end{tablalafrox}

\SIMnotario{Certifico que el presente extracto concuerda con las escrituras públicas originales de constitución y modificaciones de LafroX SpA, que he tenido a la vista. \SIMfecha. Copia autorizada simulada con fines académicos.}

\vspace{8pt}
\noindent\small Documento recreado con fines académicos. No reemplaza a la copia autorizada de las escrituras públicas que otorga un notario. Notaría, repertorios, inscripciones y personas son ficticios.
\end{document}
